\subsection{阴吹（阴道排气）：常见的"尴尬声音"，通常无害}

性生活中阴道突然排出气体、发出类似放屁的声音——俗称"阴吹"。它非常常见，多数女性一生中都会遇到，却是被谈论得最少的现象之一：很多女性第一反应是羞耻，甚至怀疑自己"下面松了"。先说结论：\textbf{这是正常的空气动力学现象，不是病，更不是"私生活不检点"的证明}。

\vspace{0.5em}
\noindent\textbf{声音从哪来}：

\begin{itemize}
	\item 性交时，阴茎的进出会把空气压入阴道穹窿（尤其后穹窿）；当体位变换、阴茎抽出或臀部抬高等使阴道内空间突然增大时，被挤压的气体冲出阴道口，振动组织发出声音——原理与放屁相似，但\textbf{气体来自外界空气，不是肠道排气}，所以没有异味
	\item \textbf{常见诱因}：某些体位（后入位、深插入、臀部抬高的体位）；分娩后盆底支撑短暂减弱；绝经后阴道壁弹性下降；个别情况与阴道较宽松有关
\end{itemize}

\vspace{0.5em}
\noindent\textbf{它不等于"阴道松弛"}：偶尔排气是几乎所有有过性生活的女性都会发生的现象，与"松弛"没有必然联系；\textbf{只有}当排气频繁、同时伴有漏尿、下坠感、性生活时明显"吸放气感"时，才提示盆底功能需要评估——处理方案是盆底肌训练（凯格尔运动），而不是忍着不说。

\vspace{0.5em}
\noindent\textbf{可以做什么}：

\begin{itemize}
	\item \textbf{调整体位与节奏}：容易"吸空气"的体位放慢节奏、减少完全抽出的动作，声响自然会减少
	\item \textbf{盆底肌训练}：规律的凯格尔运动改善阴道壁的贴合与支撑，减少气体进入（对盆底健康本身也大有好处）
	\item \textbf{心态}：把它当作"身体打了个嗝"——真正让这个瞬间变尴尬的往往不是声音本身，而是"必须假装没发生"的紧张；对它一笑而过，是最快的"治疗"
\end{itemize}

\vspace{0.5em}
\noindent\textbf{何时需要就医}：极少数情况下需要检查——排气\textbf{持续大量}且不限于性生活时；排出物带粪臭味或粪便样物质（警惕直肠阴道瘘，罕见，多与产伤、手术或放疗相关）；伴随明显的盆底松垂感、漏尿或反复阴道异常排液。传统医学文献中的"阴吹"常与脏腑虚损相联系，民间流传不少方子——若无上述红旗症状，不必为此服药。

\begin{tcolorbox}[colback=pink!5!white,colframe=red!50!black,title=贴心小叮咛]
提前知道"这是正常的"，比事后解释一百遍都有用——可以把这一节转给伴侣看：两个人都不当回事，它就真的不算事。
\end{tcolorbox}
